\section*{第 \ref{ch:g4:journey-of-food} 章——食物之旅}

\begin{solution}{exo:g4:journey-of-food:1}
口腔、\omterm{def:g4:journey-of-food:tract}{食道}、胃、\omterm{def:g4:journey-of-food:tract}{小肠}、\omterm{def:g4:journey-of-food:tract}{大肠}——
然后残渣离开身体。
\end{solution}

\begin{solution}{exo:g4:journey-of-food:2}
它由胃壁的\omterm{def:g3:bones-and-muscles:muscle}{肌肉}搅动，并与
\omterm{def:g4:journey-of-food:digestion}{消化液}混合——持续数小时。
\end{solution}

\begin{solution}{exo:g4:journey-of-food:3}
\omterm{def:g4:journey-of-food:digestion}{消化}是食物沿\omterm{def:g4:journey-of-food:tract}{消化道}转化为
足够简单、足够小、能够进入身体的成分。这些
成分就是\omterm{def:g4:journey-of-food:digestion}{营养物质}。
\end{solution}

\begin{solution}{exo:g4:journey-of-food:4}
在\omterm{def:g4:journey-of-food:tract}{小肠}中，食物通过其薄而褶皱的管壁，
进入\omterm{prop:g4:journey-of-food:absorb}{血液}。不能穿过的物质继续进入\omterm{def:g4:journey-of-food:tract}{大肠}，
在那里失去水分，并以粪便的形式排出体外。
\end{solution}

\begin{solution}{exo:g4:journey-of-food:5}
\omterm{def:g4:journey-of-food:tract}{唾液}已经把面包分解为更简单、更甜的
成分——\omterm{def:g4:journey-of-food:digestion}{消化}从口腔开始，
而甜味就是它的味道。
\end{solution}

\begin{solution}{exo:g4:journey-of-food:6}
因为进入\omterm{prop:g4:journey-of-food:absorb}{血液}需要管壁的表面积：
长达数米的肠道与深深的内褶，
使进入腹部的吸收面积变得极大——
摊平后约相当于一间小屋的地毯。
\end{solution}

\begin{solution}{exo:g4:journey-of-food:7}
例如：认真咀嚼（胃里没有牙齿）；按时、平静地进食；
（\omterm{def:g4:journey-of-food:tract}{消化道}在从容时工作得最好）；多喝水（\omterm{def:g4:journey-of-food:digestion}{消化}需要水分）；
多吃膳食纤维（它能使残渣保持移动）；
大餐后休息（胃和腿
会争夺身体的\omterm{prop:g3:balanced-diet:jobs}{能量}）。
\end{solution}

\begin{solution}{exo:g4:journey-of-food:8}
开放性答案。通常在口腔中停留的时间远少于半小时——
其余的旅程则需要一到两天。
口腔所占的时间很短，但作用决定性：
它是唯一有牙齿的环节。
\end{solution}

\begin{solution}{exo:g4:journey-of-food:9}
问题所在：没有充分咀嚼（午饭吃得太急，食物没嚼碎），
以及饭后休息不足（吃完立刻奔跑）。
改进办法：细嚼慢咽，饭后让身体平静下来再剧烈运动。
\end{solution}

\begin{solution}{exo:g4:journey-of-food:10}
咽下的食物仍在管道之内——
就像甜甜圈中间的孔，还没有真正进入身体。
它必须先被\omterm{def:g4:journey-of-food:digestion}{消化}为\omterm{def:g4:journey-of-food:digestion}{营养物质}；只有在\omterm{def:g4:journey-of-food:tract}{小肠}中，
这些\omterm{def:g4:journey-of-food:digestion}{营养物质}才会穿过管壁进入\omterm{prop:g4:journey-of-food:absorb}{血液}。
进入口腔并不等于进入身体。
\end{solution}

\begin{solution}{exo:g4:journey-of-food:11}
\omterm{def:g4:journey-of-food:digestion}{消化}必须把食物分解为\omterm{def:g4:journey-of-food:digestion}{营养物质}，而草远比面包难分解，
纤维坚韧，
每一口的营养很少
（\cref{prop:g2:what-animals-eat:daily}）。因此，食草\omterm{def:g1:animals-around-us:animal}{动物}需要更强大的“设备”：
反复使用的研磨\omterm{def:g2:teeth-and-eating-well:kinds}{臼齿}（反刍时再次咀嚼）、
用于长时间加工的四个胃室，
以及能分解自身\omterm{def:g4:journey-of-food:digestion}{消化液}无法分解之物质的肠道微\omterm{def:g1:living-or-not:living}{生物}。
食物更难处理，工厂就要更大。
\end{solution}
